\documentclass{article}

\usepackage{fancyhdr}
\usepackage{extramarks}
\usepackage{amsmath}
\usepackage{amsthm}
\usepackage{amsfonts}
\usepackage{tikz}
\usepackage[plain]{algorithm}
\usepackage{algpseudocode}

\usetikzlibrary{automata,positioning}

%
% Basic Document Settings
%

\topmargin=-0.45in
\evensidemargin=0in
\oddsidemargin=0in
\textwidth=6.5in
\textheight=9.0in
\headsep=0.25in

\linespread{1.1}

\pagestyle{fancy}
\lhead{\hmwkAuthorName}
\chead{\hmwkClass\ (\hmwkClassInstructor): \hmwkTitle}
\rhead{\firstxmark}
\lfoot{\lastxmark}
\cfoot{\thepage}

\renewcommand\headrulewidth{0.4pt}
\renewcommand\footrulewidth{0.4pt}

\setlength\parindent{0pt}

%
% Create Problem Sections
%

\newcommand{\enterProblemHeader}[1]{
	\nobreak\extramarks{}{Problem \arabic{#1} continued on next page\ldots}\nobreak{}
	\nobreak\extramarks{Problem \arabic{#1} (continued)}{Problem \arabic{#1} continued on next page\ldots}\nobreak{}
}

\newcommand{\exitProblemHeader}[1]{
	\nobreak\extramarks{Problem \arabic{#1} (continued)}{Problem \arabic{#1} continued on next page\ldots}\nobreak{}
	\stepcounter{#1}
	\nobreak\extramarks{Problem \arabic{#1}}{}\nobreak{}
}

\setcounter{secnumdepth}{0}
\newcounter{partCounter}
\newcounter{homeworkProblemCounter}
\setcounter{homeworkProblemCounter}{1}
\nobreak\extramarks{Problem \arabic{homeworkProblemCounter}}{}\nobreak{}

%
% Homework Problem Environment
%
% This environment takes an optional argument. When given, it will adjust the
% problem counter. This is useful for when the problems given for your
% assignment aren't sequential. See the last 3 problems of this template for an
% example.
%
\newenvironment{homeworkProblem}[1][-1]{
	\ifnum#1>0
		\setcounter{homeworkProblemCounter}{#1}
	\fi
	\section{Problem \arabic{homeworkProblemCounter}}
	\setcounter{partCounter}{1}
	\enterProblemHeader{homeworkProblemCounter}
}{
	\exitProblemHeader{homeworkProblemCounter}
}

%
% Homework Details
%   - Title
%   - Due date
%   - Class
%   - Section/Time
%   - Instructor
%   - Author
%

\newcommand{\hmwkTitle}{Asignacion\ \#7}
\newcommand{\hmwkDueDate}{Noviembre 14, 2024}
\newcommand{\hmwkClass}{MATE 5150}
\newcommand{\hmwkClassInstructor}{Dr. Pedro Vasquez}
\newcommand{\hmwkAuthorName}{\textbf{Alejandro Ouslan}}

%
% Title Page
%

\title{
	\vspace{2in}
	\textmd{\textbf{\hmwkClass:\ \hmwkTitle}}\\
	\normalsize\vspace{0.1in}\small{Due\ on\ \hmwkDueDate}\\
	\vspace{0.1in}\large{\textit{\hmwkClassInstructor}}
	\vspace{3in}
}

\author{\hmwkAuthorName}
\date{}

\renewcommand{\part}[1]{\textbf{\large Part \Alph{partCounter}}\stepcounter{partCounter}\\}

%
% Various Helper Commands
%

% Useful for algorithms
\newcommand{\alg}[1]{\textsc{\bfseries \footnotesize #1}}

% For derivatives
\newcommand{\deriv}[1]{\frac{\mathrm{d}}{\mathrm{d}x} (#1)}

% For partial derivatives
\newcommand{\pderiv}[2]{\frac{\partial}{\partial #1} (#2)}

% Integral dx
\newcommand{\dx}{\mathrm{d}x}

% Alias for the Solution section header
\newcommand{\solution}{\textbf{\large Solution}}

% Probability commands: Expectation, Variance, Covariance, Bias
\newcommand{\E}{\mathrm{E}}
\newcommand{\Var}{\mathrm{Var}}
\newcommand{\Cov}{\mathrm{Cov}}
\newcommand{\Bias}{\mathrm{Bias}}

\begin{document}

\maketitle

\pagebreak

% Homework problem 1 sec 5.1 problem 3c
\begin{homeworkProblem}
	For each of the following linear operators $T$ on a vector space $V$ and ordered basis $\beta$, compute $[T]_{\beta}$ and determine
	whether $\beta$ is a basis consisting of eigenvectors of $T$.
	\[
		V = \mathbb{R}^3, T \begin{pmatrix} a \\ b \\ c\end{pmatrix} =
		\begin{pmatrix} 3a + 2b - 2c \\ -4a - 3b + 2c \\ -c \end{pmatrix}, \text{ and }
		\beta = \left\{ \begin{pmatrix} 0 \\ 1 \\ 1 \end{pmatrix}, \begin{pmatrix} 1 \\ -1 \\ 0 \end{pmatrix}, \begin{pmatrix} 1 \\ 0 \\ 2 \end{pmatrix} \right\}
	\]
	Compute $[T]_{\beta}$:
	\[
		\begin{split}
			\begin{pmatrix} 3(0) + 2(1) - 2(1) \\ -4(0) - 3(1) + 2(1) \\ -(1) \end{pmatrix}                                                       & = \begin{pmatrix} 0 \\ -1 \\ -1 \end{pmatrix}                         \\
			a\begin{pmatrix} 0 \\ 1 \\ 1 \end{pmatrix} + b\begin{pmatrix} 1 \\ -1 \\ 0 \end{pmatrix} + c\begin{pmatrix} 1 \\ 0 \\ 2 \end{pmatrix} & = \begin{pmatrix} 0 \\ -1 \\ -1 \end{pmatrix}                         \\
			v_1                                                                                                                                   & = \begin{pmatrix} -1 \\ 0 \\ 1 \end{pmatrix}                          \\
			v_2                                                                                                                                   & = \begin{pmatrix} 0 \\ 1 \\ 0 \end{pmatrix}                           \\
			v_3                                                                                                                                   & = \begin{pmatrix} 0 \\ 0 \\ -1 \end{pmatrix}                          \\
			[T]_\beta                                                                                                                             & = \begin{pmatrix} -1 & 0 & 0 \\ 0 & 1 & 0 \\ 0 & 0 & -1 \end{pmatrix} \\
		\end{split}
	\]
	Finding the eigenvalues of $T$:
	\[
		\begin{split}
			\text{det}([T]_\beta - \lambda I) & = \text{det}\left(\begin{pmatrix} -1 - \lambda & 0 & 0 \\ 0 & 1 - \lambda & 0 \\ 0 & 0 & -1 - \lambda \end{pmatrix}\right) \\
			                                  & = (1-\lambda)(1+\lambda)^2 = 0                                                                                             \\
			\lambda_1, \lambda_2              & = 1, -1
		\end{split}
	\]
	Calculating the eigenvectors for $\lambda_1 = 1$:
	\[
		\begin{split}
			\begin{pmatrix} -1 - 1 & 0 & 0 \\ 0 & 1 - 1 & 0 \\ 0 & 0 & -1 - 1 \end{pmatrix}\begin{pmatrix} x \\ y \\ z \end{pmatrix} & = \begin{pmatrix} 0 \\ 0 \\ 0 \end{pmatrix} \\
			\begin{pmatrix} -2 & 0 & 0 \\ 0 & 0 & 0 \\ 0 & 0 & -2 \end{pmatrix}\begin{pmatrix} x \\ y \\ z \end{pmatrix}             & = \begin{pmatrix} 0 \\ 0 \\ 0 \end{pmatrix} \\
			v_1                                                                                                                      & = \begin{pmatrix} 0 \\ 1 \\ 0 \end{pmatrix}
		\end{split}
	\]
	Calculating the eigenvectors for $\lambda_2 = -1$:
	\[
		\begin{split}
			\begin{pmatrix} -1 + 1 & 0 & 0 \\ 0 & 1 + 1 & 0 \\ 0 & 0 & -1 + 1 \end{pmatrix}\begin{pmatrix} x \\ y \\ z \end{pmatrix} & = \begin{pmatrix} 0 \\ 0 \\ 0 \end{pmatrix} \\
			\begin{pmatrix} 0 & 0 & 0 \\ 0 & 2 & 0 \\ 0 & 0 & 0 \end{pmatrix}\begin{pmatrix} x \\ y \\ z \end{pmatrix}               & = \begin{pmatrix} 0 \\ 0 \\ 0 \end{pmatrix} \\
			v_2                                                                                                                      & = \begin{pmatrix} 1 \\ 0 \\ 0 \end{pmatrix} \\
			v_3                                                                                                                      & = \begin{pmatrix} 0 \\ 0 \\ 1 \end{pmatrix}
		\end{split}
	\]

\end{homeworkProblem}

% Homework problem 2 sec 5.1 problem 4b
\begin{homeworkProblem}
	For each of the following matrices $A \in M_{n \times n}(F)$,
	\begin{enumerate}
		\item Determine all the eigenvalues of $A$.
		\item For each eigenvalue $\lambda$ of $A$, find the set of eigenvecotrs corresponding to $\lambda$.
		\item If possible, find a basis for $F^n$ consisting of eigenvectors of $A$.
		\item If successful in finding such a basis, determine an invertible matrix $Q$ and a diagonal matrix $D$ such that $Q^{-1}AQ = D$.
	\end{enumerate}
	\[
		A = \begin{pmatrix} 0 & -2 & -3 \\ -1 & 1 & -1 \\ 2 & 2 & 5 \end{pmatrix} \text{ for } F = \mathbb{R}
	\]
	Determine all the eigenvalues of $A$:
	\[
		\begin{split}
			\text{det}(A - \lambda I)       & = \text{det}\left(\begin{pmatrix} 0 - \lambda & -2 & -3 \\ -1 & 1 - \lambda & -1 \\ 2 & 2 & 5 - \lambda \end{pmatrix}\right) \\
			                                & = (3-\lambda)(2-\lambda)(1-\lambda) = 0                                                                                      \\
			\lambda_1, \lambda_2, \lambda_3 & = 3, 2, 1
		\end{split}
	\]
	Calculating the eigenvectors for $\lambda_1 = 3$:
	\[
		\begin{split}
			\begin{pmatrix} 0 - 3 & -2 & -3 \\ -1 & 1 - 3 & -1 \\ 2 & 2 & 5 - 3 \end{pmatrix}\begin{pmatrix} x \\ y \\ z \end{pmatrix} & = \begin{pmatrix} 0 \\ 0 \\ 0 \end{pmatrix}  \\
			\begin{pmatrix} -3 & -2 & -3 \\ -1 & -2 & -1 \\ 2 & 2 & 2 \end{pmatrix}\begin{pmatrix} x \\ y \\ z \end{pmatrix}           & = \begin{pmatrix} 0 \\ 0 \\ 0 \end{pmatrix}  \\
			\begin{pmatrix} -3 & -2 & -3 \\ 0 & 4 & 0 \\ 0 & 0 & 0 \end{pmatrix}\begin{pmatrix} x \\ y \\ z \end{pmatrix}              & = \begin{pmatrix} 0 \\ 0 \\ 0 \end{pmatrix}  \\
			v_1                                                                                                                        & = \begin{pmatrix} -1 \\ 0 \\ 1 \end{pmatrix} \\
		\end{split}
	\]
	Calculating the eigenvectors for $\lambda_2 = 2$:
	\[
		\begin{split}
			\begin{pmatrix} 0 - 2 & -2 & -3 \\ -1 & 1 - 2 & -1 \\ 2 & 2 & 5 - 2 \end{pmatrix}\begin{pmatrix} x \\ y \\ z \end{pmatrix} & = \begin{pmatrix} 0 \\ 0 \\ 0 \end{pmatrix}  \\
			\begin{pmatrix} -2 & -2 & -3 \\ -1 & -1 & -1 \\ 2 & 2 & 3 \end{pmatrix}\begin{pmatrix} x \\ y \\ z \end{pmatrix}           & = \begin{pmatrix} 0 \\ 0 \\ 0 \end{pmatrix}  \\
			\begin{pmatrix} -2 & -2 & -3 \\ 0 & 0 & -1 \\ 0 & 0 & 0 \end{pmatrix}\begin{pmatrix} x \\ y \\ z \end{pmatrix}             & = \begin{pmatrix} 0 \\ 0 \\ 0 \end{pmatrix}  \\
			v_2                                                                                                                        & = \begin{pmatrix} -1 \\ 1 \\ 0 \end{pmatrix}
		\end{split}
	\]
	Calculating the eigenvectors for $\lambda_3 = 1$:
	\[
		\begin{split}
			\begin{pmatrix} 0 - 1 & -2 & -3 \\ -1 & 1 - 1 & -1 \\ 2 & 2 & 5 - 1 \end{pmatrix}\begin{pmatrix} x \\ y \\ z \end{pmatrix} & = \begin{pmatrix} 0 \\ 0 \\ 0 \end{pmatrix}  \\
			\begin{pmatrix} -1 & -2 & -3 \\ -1 & 0 & -1 \\ 2 & 2 & 4 \end{pmatrix}\begin{pmatrix} x \\ y \\ z \end{pmatrix}            & = \begin{pmatrix} 0 \\ 0 \\ 0 \end{pmatrix}  \\
			\begin{pmatrix} -2 & -2 & -3 \\ 0 & 0 & -1 \\ 0 & 0 & 0 \end{pmatrix}\begin{pmatrix} x \\ y \\ z \end{pmatrix}             & = \begin{pmatrix} 0 \\ 0 \\ 0 \end{pmatrix}  \\
			v_3                                                                                                                        & = \begin{pmatrix} -1 \\ 1 \\ 1 \end{pmatrix}
		\end{split}
	\]
	The basis of eigenvectors for $A$ is $\left\{ \begin{pmatrix} -1 \\ 0 \\ 1 \end{pmatrix}, \begin{pmatrix} -1 \\ 1 \\ 0 \end{pmatrix}, \begin{pmatrix} -1 \\ 1 \\ 1 \end{pmatrix} \right\}$.
	\[
		\begin{split}
			Q        & = \begin{pmatrix} -1 & -1 & -1 \\ 0 & 1 & 1 \\ 1 & 0 & 1 \end{pmatrix}   \\
			Q^{-1}   & = \begin{pmatrix} -1 & -1 & -1 \\ -1 & 0 & -1 \\ 1 & 1 & 2 \end{pmatrix} \\
			Q^{-1}AQ & = \begin{pmatrix} 1 & 0 & 0 \\ 0 & 2 & 0 \\ 0 & 0 & 3 \end{pmatrix}
		\end{split}
	\]
\end{homeworkProblem}

% Homework problem 3 sec 5.1 problem 5i
\begin{homeworkProblem}
	For each linear operator $T$ on $V$, find the eigenvalues of $T$ and an ordered basis $\beta$ for $V$ such that $[T]_{\beta}$ is a diagonal matrix.
	\[
		V = M_{2 \times 2}(\mathbb{R}), \text{ and } T \begin{pmatrix} a & b \\ c & d \end{pmatrix} = \begin{pmatrix} c & d \\ a & b \end{pmatrix}
	\]
	\[
		\begin{split}
			\beta                                        & = \left \{ \begin{pmatrix} 1 & 0 \\ 0 & 0 \end{pmatrix}, \begin{pmatrix} 0 & 1 \\ 0 & 0 \end{pmatrix}, \begin{pmatrix} 0 & 0 \\ 1 & 0 \end{pmatrix}, \begin{pmatrix} 0 & 0 \\ 0 & 1 \end{pmatrix} \right \}         \\
			\begin{pmatrix} 1 & 0 \\ 0 & 0 \end{pmatrix} & = x_1 \begin{pmatrix} 1 & 0 \\ 0 & 0 \end{pmatrix} + x_2 \begin{pmatrix} 0 & 1 \\ 0 & 0 \end{pmatrix} + x_3 \begin{pmatrix} 0 & 0 \\ 1 & 0 \end{pmatrix} + x_4 \begin{pmatrix} 0 & 0 \\ 0 & 1 \end{pmatrix}         \\
			\beta                                        & = \left \{ \begin{pmatrix} 0 \\ 0 \\ 1 \\ 0 \end{pmatrix}, \begin{pmatrix} 0 \\ 0 \\ 0 \\ 1 \end{pmatrix}, \begin{pmatrix} 1 \\ 0 \\ 0 \\ 0 \end{pmatrix}, \begin{pmatrix} 0 \\ 1 \\ 0 \\ 0 \end{pmatrix} \right \} \\
			[T]_\beta                                    & = \begin{pmatrix} 0 & 0 & 1 & 0 \\ 0 & 0 & 0 & 1 \\ 1 & 0 & 0 & 0 \\ 0 & 1 & 0 & 0 \end{pmatrix}
		\end{split}
	\]
	Calculating the eigenvalues of $T$:
	\[
		\begin{split}
			\text{det}(T - \lambda I) & = \text{det}\left(\begin{pmatrix} 0 - \lambda & 0 & 1 & 0 \\ 0 & 0 - \lambda & 0 & 1 \\ 1 & 0 & 0 - \lambda & 0 \\ 0 & 1 & 0 & 0 - \lambda \end{pmatrix}\right) \\
			                          & = (1 - \lambda)^2(\lambda - 1)^2 = 0                                                                                                                            \\
		\end{split}
	\]
	Calculating the eigenvectors for $\lambda_1 = 1$:
	\[
		\begin{split}
			\begin{pmatrix} 0 - 1 & 0 & 1 & 0 \\ 0 & 0 - 1 & 0 & 1 \\ 1 & 0 & 0 - 1 & 0 \\ 0 & 1 & 0 & 0 - 1 \end{pmatrix}\begin{pmatrix} x \\ y \\ z \\ w \end{pmatrix} & = \begin{pmatrix} 0 \\ 0 \\ 0 \\ 0 \end{pmatrix} \\
			\begin{pmatrix} -1 & 0 & 1 & 0 \\ 0 & -1 & 0 & 1 \\ 1 & 0 & -1 & 0 \\ 0 & 1 & 0 & -1 \end{pmatrix}\begin{pmatrix} x \\ y \\ z \\ w \end{pmatrix}             & = \begin{pmatrix} 0 \\ 0 \\ 0 \\ 0 \end{pmatrix} \\
			v = a\begin{pmatrix} 1 \\ 0 \\ 1 \\ 0 \end{pmatrix} + b\begin{pmatrix} 0 \\ 1 \\ 0 \\ 1 \end{pmatrix}
		\end{split}
	\]
	Calculating the eigenvectors for $\lambda_2 = 1$:
	\[
		\begin{split}
			\begin{pmatrix} 0 - 1 & 0 & 1 & 0 \\ 0 & 0 - 1 & 0 & 1 \\ 1 & 0 & 0 - 1 & 0 \\ 0 & 1 & 0 & 0 - 1 \end{pmatrix}\begin{pmatrix} x \\ y \\ z \\ w \end{pmatrix} & = \begin{pmatrix} 0 \\ 0 \\ 0 \\ 0 \end{pmatrix} \\
			\begin{pmatrix} -1 & 0 & 1 & 0 \\ 0 & -1 & 0 & 1 \\ 1 & 0 & -1 & 0 \\ 0 & 1 & 0 & -1 \end{pmatrix}\begin{pmatrix} x \\ y \\ z \\ w \end{pmatrix}             & = \begin{pmatrix} 0 \\ 0 \\ 0 \\ 0 \end{pmatrix} \\
			v = a\begin{pmatrix} -1 \\ 0 \\ 1 \\ 0 \end{pmatrix} + b\begin{pmatrix} 0 \\ -1 \\ 0 \\ 1 \end{pmatrix}
		\end{split}
	\]
	V is a basis of eigenvectors for $T$. $V = \begin{pmatrix} 1 & 0 & 0 & 0 \\ 0 & 1 & 0 & 0 \\ 0 & 0 & 1 & 0 \\ 0 & 0 & 0 & 1 \end{pmatrix}$.

\end{homeworkProblem}

% Homework problem 4 sec 5.1 problem 11
\begin{homeworkProblem}
	Let $V$ be a finite-dimensional vector space, and let $\lambda$ be any scalar.
	\begin{enumerate}
		\item For any ordered basis $\beta$ for $V$, prove that $[\lambda|_{V}]_{\beta} = \lambda I$. \\
		      \textbf{Solution:}
		      Consider the linear map $\lambda|_V$ defined by $\lambda|_V(v) = \lambda v$ for all $v \in V$. Let $\beta = \{v_1, v_2, \dots, v_n\}$ be an ordered basis of $V$.
		      For each $v_i \in \beta$, we have
		      \[
			      \lambda|_V(v_i) = \lambda v_i.
		      \]
		      In terms of the basis $\beta$, this corresponds to the vector
		      \[
			      [\lambda|_V(v_i)]_{\beta} = \lambda [v_i]_{\beta} = \lambda e_i,
		      \]
		      where $e_i$ is the $i$-th standard basis vector.
		      Therefore, the matrix representation of $\lambda|_V$ in the basis $\beta$ is
		      \[
			      [\lambda|_V]_{\beta} = \lambda I,
		      \]
		      where $I$ is the identity matrix of size $n \times n$.
		\item Compute the characteristic polynomial of $\lambda|_V$.
		      \textbf{Solution:}
		      The characteristic polynomial of a linear operator $T$ is given by
		      \[
			      p_T(\mu) = \det(\mu I - [T]_{\beta}).
		      \]
		      In this case, for the operator $\lambda|_V$, we have
		      \[
			      p_{\lambda|_V}(\mu) = \det(\mu I - \lambda I) = \det((\mu - \lambda)I) = (\mu - \lambda)^n.
		      \]
		      Thus, the characteristic polynomial of $\lambda|_V$ is
		      \[
			      p_{\lambda|_V}(\mu) = (\mu - \lambda)^n.
		      \]
		\item Show that $\lambda|_V$ is diagonalizable and has only one eigenvalue. \\
		      \textbf{Solution:}
		      From the characteristic polynomial, we see that the only eigenvalue of $\lambda|_V$ is $\lambda$, with algebraic multiplicity $n$.
		      Since the matrix representation of $\lambda|_V$ in the basis $\beta$ is $\lambda I$, which is already a diagonal matrix, $\lambda|_V$ is diagonalizable.
		      Therefore, $\lambda|_V$ is diagonalizable and has only one eigenvalue, which is $\lambda$.
	\end{enumerate}
\end{homeworkProblem}

% Homework problem 5 sec 5.1 problem 22
\begin{homeworkProblem}
	Let $T$ be a linear operator on a vector space $V$ over the field $F$, and let $g(t)$ be a polynomial with coefficients from $F$. Prove that if $x$ is an eigenvector of $T$ with corresponding eigenvalue $\lambda$, then $g(T)(x) = g(\lambda)x$. That is, $x$ is an eigenvector of $g(T)$ with corresponding eigenvalue $g(\lambda)$. \\
	\textbf{Solution:}
	Let the polynomial $g(t)$ be given by
	\[
		g(t) = a_n t^n + a_{n-1} t^{n-1} + \cdots + a_1 t + a_0,
	\]
	where the coefficients \( a_i \) belong to the field \( F \).
	Since \( x \) is an eigenvector of \( T \) with eigenvalue \( \lambda \), we have
	\[
		T(x) = \lambda x.
	\]
	We want to prove that \( g(T)(x) = g(\lambda)x \). First, observe that
	\[
		g(T) = a_n T^n + a_{n-1} T^{n-1} + \cdots + a_1 T + a_0 I,
	\]
	where \( I \) is the identity operator on \( V \). Applying \( g(T) \) to \( x \), we get
	\[
		g(T)(x) = a_n T^n(x) + a_{n-1} T^{n-1}(x) + \cdots + a_1 T(x) + a_0 x.
	\]
	Using the fact that \( T(x) = \lambda x \), we compute the higher powers of \( T \) acting on \( x \):
	\[
		T^k(x) = \lambda^k x \quad \text{for any integer } k \geq 0.
	\]
	Therefore, each term in the sum becomes
	\[
		a_n T^n(x) = a_n \lambda^n x, \quad a_{n-1} T^{n-1}(x) = a_{n-1} \lambda^{n-1} x, \quad \dots, \quad a_1 T(x) = a_1 \lambda x, \quad a_0 x = a_0 x.
	\]
	Combining all these terms, we obtain
	\[
		g(T)(x) = (a_n \lambda^n + a_{n-1} \lambda^{n-1} + \cdots + a_1 \lambda + a_0) x = g(\lambda) x.
	\]
	Thus, we have shown that
	\[
		g(T)(x) = g(\lambda) x.
	\]
	Therefore, \( x \) is an eigenvector of \( g(T) \) with corresponding eigenvalue \( g(\lambda) \).
\end{homeworkProblem}

% Homework problem 6 sec 5.2 problem 2g
\begin{homeworkProblem}
	For each of the following matrices $A \in M_{n \times n}(\mathbb{R})$, test $A$ for diagonalizability, and if $A$ is diagonalizable, find an invertible matrix
	$Q$ and a diagonal matrix $D$ such that $Q^{-1}AQ = D$.
	\[
		A = \begin{pmatrix} 3 & 1 & 1 \\ 2 & 4 & 2 \\ -1 & -1 & 1 \end{pmatrix}
	\]
\end{homeworkProblem}

% Homework problem 7 sec 5.2 problem 7
\begin{homeworkProblem}
	For
	\[
		A = \begin{pmatrix} 1 & 4 \\ 2 & 3 \end{pmatrix} \in M_{2 \times 2}(\mathbb{R}),
	\]
	find a expression for $A^n$, where $n$ is an arbitrary positive integer.
\end{homeworkProblem}

% Homework problem 8 sec 5.2 problem 9
\begin{homeworkProblem}
	Let $T$ be a linear operator on a finite-dimensional vector space $V$, and suppose there exists an ordered basis $\beta$ for $V$ such that $[T]_{\beta}$ is
	an upper triangular matrix.
	\begin{enumerate}
		\item Prove that the caracteristic polynomial for $T$ splits.
		\item State and prove an analogous result for matrices.
	\end{enumerate}
	The converse of (a) is treated in Exercise 12(b).
\end{homeworkProblem}

% Homework problem 9 sec 5.4 problem 2d
\begin{homeworkProblem}
	For each of the following linear operators $T$ on the vector space $V$, determine whether the given subspace $W$ is a $T$-invariant subspace of $V$.
	\[
		V = C([0, 1]), T(f(t)) = \left[\int_{0}^{1} f(x)dt\right]t, \text{ and } W = \{f \in V : f(t) = at + b \text{ for some } a, b \}
	\]
\end{homeworkProblem}

% Homework problem 10 sec 5.4 problem 6d
\begin{homeworkProblem}
	For each linear operator $T$ on the vector space $V$, find an ordered basis for the $T$-cyclic subspace generated by the given vector $z$
	\[
		V = M_{2 \times 2}(\mathbb{R}), T(A) = \begin{pmatrix} 1 & 1 \\ 2 & 2 \end{pmatrix}A, \text{ and } z = \begin{pmatrix} 0 & 1 \\ 1 & 0 \end{pmatrix}
	\]
\end{homeworkProblem}

% Homework problem 11 sec 5.4 problem 10d
\begin{homeworkProblem}
	For each linear operator in Exercise 6, find the characteristic polynomial of $f(t)$ of $T$, and verify that the characteristic polynomial of $T_{W}$
	(computed in Exercise 6) divides $f(t)$.
\end{homeworkProblem}

% Homework problem 12 sec 5.4 problem 11
\begin{homeworkProblem}
	Let \( T \) be a linear operator on a vector space \( V \), let \( v \) be a nonzero vector in \( V \), and let \( W \) be the \( T \)-cyclic subspace of \( V \) generated by \( v \). Prove that
	\begin{enumerate}
		\item \( W \) is \( T \)-invariant.
		\item Any \( T \)-invariant subspace of \( V \) containing \( v \) also contains \( W \).
	\end{enumerate}
	\textbf{Solution:}
	\textbf{(1) Prove that \( W \) is \( T \)-invariant.}
	The \( T \)-cyclic subspace \( W \) generated by \( v \) is the span of all the vectors of the form \( T^k v \), where \( k \geq 0 \). That is,
	\[
		W = \text{span} \{ v, T v, T^2 v, \dots \}.
	\]
	We need to show that \( W \) is \( T \)-invariant, meaning that \( T(W) \subseteq W \). This means that for any vector \( w \in W \), we must show that \( T(w) \in W \).
	Let \( w \in W \). By the definition of \( W \), \( w \) is a linear combination of the vectors \( v, T v, T^2 v, \dots \), i.e.,
	\[
		w = a_0 v + a_1 T v + a_2 T^2 v + \dots + a_n T^n v,
	\]
	where \( a_0, a_1, \dots, a_n \) are scalars.
	Applying \( T \) to \( w \), we get
	\[
		T(w) = a_0 T(v) + a_1 T^2(v) + a_2 T^3(v) + \dots + a_n T^{n+1} v.
	\]
	Since each of the terms \( T^{k+1} v \) (for \( k = 0, 1, \dots, n \)) is still a linear combination of the vectors \( v, T v, T^2 v, \dots \), it follows that \( T(w) \) is also a linear combination of the vectors \( v, T v, T^2 v, \dots \), which are in \( W \). Therefore, \( T(w) \in W \), and we conclude that \( W \) is \( T \)-invariant.
	\textbf{(2) Prove that any \( T \)-invariant subspace of \( V \) containing \( v \) also contains \( W \).}
	Let \( U \) be any \( T \)-invariant subspace of \( V \) that contains \( v \). We want to show that \( U \) contains \( W \).
	By the definition of \( W \), we know that
	\[
		W = \text{span} \{ v, T v, T^2 v, \dots \}.
	\]
	Since \( U \) is \( T \)-invariant and contains \( v \), it must also contain \( T v \), \( T^2 v \), and all higher powers of \( T \) applied to \( v \). Specifically, since \( U \) is closed under \( T \), we have
	\[
		T v \in U, \quad T^2 v \in U, \quad \dots.
	\]
	Thus, every vector in the span of \( \{ v, T v, T^2 v, \dots \} \), which defines \( W \), must lie in \( U \). Hence, \( W \subseteq U \).
	Therefore, any \( T \)-invariant subspace of \( V \) containing \( v \) also contains \( W \).
\end{homeworkProblem}

% Homework problem 13 sec 6.1 problem 3
\begin{homeworkProblem}
	In $C([0, 1])$, let $f(t) = t$ and $g(t) = e^{t}$. Compute $\langle f, g \rangle$. (as defined in Example 3), $\|f\|$, $\|g\|$, and $\|f + g\|$. Then verify
	both the Cauchy-Schwarz inequality and the triangle inequality.
\end{homeworkProblem}

% Homework problem 14 sec 6.1 problem 5
\begin{homeworkProblem}
	In $C^2$, show that $\langle x, y \rangle = xAy^{\star}$ is an inner product, where
	\[
		A = \begin{pmatrix} 1 & i \\ -i & 2 \end{pmatrix}
	\]
	Compute $\langle x, y \rangle$ for $x = (1 - i, 2 + 3i)$ and $y = (2 + i, 3 -2i)$.
\end{homeworkProblem}

% Homework problem 15 sec 6.1 problem 8b
\begin{homeworkProblem}
	Provide reasons why each of the following is not an inner product on the given vector space.
	\[
		\langle A, B \rangle + \text{tr}(A + B) \text{ on } M_{2 \times 2}(\mathbb{R}).
	\]
	\textbf{Solution:}
	An inner product on a vector space must satisfy the following properties:
	1. **Linearity in the first argument**: For all \( A, B, C \in M_{2 \times 2}(\mathbb{R}) \) and scalars \( \alpha \in \mathbb{R} \),
	\[
		\langle \alpha A + B, C \rangle = \alpha \langle A, C \rangle + \langle B, C \rangle.
	\]
	2. **Symmetry**: For all \( A, B \in M_{2 \times 2}(\mathbb{R}) \),
	\[
		\langle A, B \rangle = \langle B, A \rangle.
	\]

	3. **Positive definiteness**: For all \( A \in M_{2 \times 2}(\mathbb{R}) \),
	\[
		\langle A, A \rangle \geq 0 \quad \text{with equality if and only if } A = 0.
	\]

	Let's examine why the given expression \( \langle A, B \rangle + \text{tr}(A + B) \) is **not** an inner product.

	**1. Violation of Positive Definiteness:**

	Consider the expression for the inner product:
	\[
		\langle A, B \rangle + \text{tr}(A + B).
	\]
	To be a valid inner product, it must satisfy positive definiteness, meaning that
	\[
		\langle A, A \rangle + \text{tr}(A + A) \geq 0,
	\]
	with equality if and only if \( A = 0 \).

	However, observe the second term, \( \text{tr}(A + A) = \text{tr}(2A) = 2 \text{tr}(A) \). This term can be nonzero even when \( A \neq 0 \), as \( \text{tr}(A) \) is generally not zero for a nonzero matrix \( A \). For example, if \( A = \begin{pmatrix} 1 & 0 \\ 0 & 0 \end{pmatrix} \), then \( \text{tr}(A) = 1 \), and hence,
	\[
		\langle A, A \rangle + \text{tr}(A + A) = \langle A, A \rangle + 2 \cdot 1 \neq 0.
	\]
	This violates the positive definiteness requirement, because for an inner product, the value should be zero only when \( A = 0 \), but in this case, it is nonzero for nonzero matrices.

	**2. Violation of Linearity in the First Argument:**

	The expression \( \langle A, B \rangle + \text{tr}(A + B) \) does not satisfy linearity in the first argument. Specifically, consider the second term \( \text{tr}(A + B) \). For a scalar \( \alpha \), the term \( \text{tr}(\alpha A + B) \) is not equal to \( \alpha \text{tr}(A + B) \), as the trace of a sum is linear in each individual matrix, but this term introduces an additive constant that doesn't preserve linearity in both arguments simultaneously.

	For example, for \( A = \begin{pmatrix} 1 & 0 \\ 0 & 0 \end{pmatrix} \), \( B = \begin{pmatrix} 0 & 0 \\ 0 & 1 \end{pmatrix} \), and \( \alpha = 2 \), we have:
	\[
		\text{tr}(A + B) = \text{tr}\left( \begin{pmatrix} 1 & 0 \\ 0 & 1 \end{pmatrix} \right) = 2,
	\]
	but
	\[
		\text{tr}(\alpha A + B) = \text{tr}\left( \begin{pmatrix} 2 & 0 \\ 0 & 0 \end{pmatrix} + \begin{pmatrix} 0 & 0 \\ 0 & 1 \end{pmatrix} \right) = \text{tr}\left( \begin{pmatrix} 2 & 0 \\ 0 & 1 \end{pmatrix} \right) = 3.
	\]
	Hence, the term \( \text{tr}(A + B) \) is not linear in the first argument, violating the linearity condition for inner products.

	**Conclusion:**

	The expression \( \langle A, B \rangle + \text{tr}(A + B) \) fails to satisfy the positive definiteness property and the linearity property in the first argument. Therefore, it is not an inner product on \( M_{2 \times 2}(\mathbb{R}) \).
\end{homeworkProblem}

% Homework problem 16 sec 6.1 problem 13
\begin{homeworkProblem}
	Suppose that $\langle \cdot, \cdot \rangle_1$ and $\langle \cdot, \cdot \rangle_2$ are two inner products on a vector space $V$. Prove that
	$\langle \cdot, \cdot \rangle = \langle \cdot, \cdot \rangle_1 + \langle \cdot, \cdot \rangle_2$ is another inner product on $V$.
	\begin{proof}
		Let $\langle \cdot, \cdot \rangle = \langle \cdot, \cdot \rangle_1 + \langle \cdot, \cdot \rangle_2$. We need to show that $\langle \cdot, \cdot \rangle$ satisfies the properties of an inner product:

		1. **Linearity in the first argument:**

		For any $u, v, w \in V$ and scalar $\alpha$, we have:
		\[
			\langle u + \alpha v, w \rangle = \langle u + \alpha v, w \rangle_1 + \langle u + \alpha v, w \rangle_2
		\]
		By the linearity of $\langle \cdot, \cdot \rangle_1$ and $\langle \cdot, \cdot \rangle_2$:
		\[
			= \left( \langle u, w \rangle_1 + \alpha \langle v, w \rangle_1 \right) + \left( \langle u, w \rangle_2 + \alpha \langle v, w \rangle_2 \right)
		\]
		\[
			= \langle u, w \rangle + \alpha \langle v, w \rangle
		\]
		Hence, $\langle \cdot, \cdot \rangle$ is linear in the first argument.

		2. **Conjugate symmetry:**

		For any $u, v \in V$:
		\[
			\langle u, v \rangle = \langle u, v \rangle_1 + \langle u, v \rangle_2
		\]
		By the conjugate symmetry of $\langle \cdot, \cdot \rangle_1$ and $\langle \cdot, \cdot \rangle_2$:
		\[
			= \overline{\langle v, u \rangle_1} + \overline{\langle v, u \rangle_2}
		\]
		\[
			= \overline{\langle v, u \rangle}
		\]
		Hence, $\langle \cdot, \cdot \rangle$ satisfies conjugate symmetry.

		3. **Positive definiteness:**

		For any $u \in V$, we have:
		\[
			\langle u, u \rangle = \langle u, u \rangle_1 + \langle u, u \rangle_2
		\]
		By the positive definiteness of $\langle \cdot, \cdot \rangle_1$ and $\langle \cdot, \cdot \rangle_2$:
		\[
			= \|u\|_1^2 + \|u\|_2^2
		\]
		Since both $\|u\|_1^2 \geq 0$ and $\|u\|_2^2 \geq 0$, we conclude that:
		\[
			\langle u, u \rangle \geq 0
		\]
		Furthermore, if $\langle u, u \rangle = 0$, then both $\|u\|_1^2 = 0$ and $\|u\|_2^2 = 0$, so $u = 0$.

		Thus, $\langle \cdot, \cdot \rangle$ satisfies all the properties of an inner product. Therefore, $\langle \cdot, \cdot \rangle$ is an inner product on $V$.
	\end{proof}
\end{homeworkProblem}

% Homework problem 17 sec 6.1 problem 17
\begin{homeworkProblem}
	Let $T$ be a linear operator on an inner product space $V$, and suppose $\|T(x)\| = \|x\|$ for all $x$. Prove that $T$ is one-to-one.
	\begin{proof}
		We are given that \( \|T(x)\| = \|x\| \) for all \( x \in V \). To prove that \( T \) is one-to-one, we need to show that if \( T(x) = T(y) \), then \( x = y \).

		Let \( x, y \in V \) and assume \( T(x) = T(y) \). This implies that \( T(x - y) = T(x) - T(y) = 0 \), so \( T(x - y) = 0 \).

		Now, we will use the given condition that \( \|T(x)\| = \|x\| \) for all \( x \in V \). Specifically, we have:
		\[
			\|T(x - y)\| = \|x - y\| \quad \text{(since \( T(x - y) = 0 \))}.
		\]
		Thus,
		\[
			\|0\| = \|x - y\| \quad \Rightarrow \quad \|x - y\| = 0.
		\]
		Since the norm of a vector is zero if and only if the vector itself is the zero vector, we conclude that \( x - y = 0 \), or equivalently, \( x = y \).

		Therefore, \( T \) is one-to-one.
	\end{proof}
\end{homeworkProblem}

% Homework problem 18 sec 6.1 problem 26c
\begin{homeworkProblem}
	Prove that the following are norms on the given vector space $V$.
	\[
		V = C([0,1]), \|f\|_{V} = \int_{0}^{1} |f(t)|dt \text{ for all } f \in V
	\]
\end{homeworkProblem}

% Homework problem 19 sec 6.2 problem 2c
\begin{homeworkProblem}
	In each part, apply the Gram-Schmidt process to the given subset $S$ of the inner product space $V$ to obtain an orthogonal basis for
	$\text{span}(S)$. Then normalize the vectors in this basis to obtain an orthonormal basis $\beta$ for $\text{span}(S)$, and compute
	the Fourier coefficients of the given vector relative to $\beta$. Finally, use Theorem 6.5 to verify your results.
	\[
		V = P_2(\mathbb{R}), \text{ with the inner product } \langle f(x), g(x) \rangle = \int_{0}^{1} f(t)g(t)dt,
		S = \{1, x, x^2\}, \text{ and } h(x) = 1 + x
	\]
\end{homeworkProblem}

% Homework problem 20 sec 6.2 problem 2i
\begin{homeworkProblem}
	Same as the previous problem
	\[
		V + span(S) \text{ with the inner product } \langle f(x), g(x) \rangle = \int_{0}^{\pi} f(t)g(t)dt,
		S = \{\sin(t), \cos(t), 1, t\}, \text{ and } h(x) = 2t + 1
	\]
\end{homeworkProblem}

\begin{homeworkProblem}
	Let $S = \{(1,0,i), (1,2,1)\}$ in $\mathbb{C}^3$. Compute $S^{\perp}$.

	\textbf{Solution:}

	We need to find the orthogonal complement of the set $S$ in $\mathbb{C}^3$. Let $v = (x_1, x_2, x_3) \in S^{\perp}$. Then, for $v$ to be in the orthogonal complement of $S$, we require:

	\[
		\langle v, (1, 0, i) \rangle = 0 \quad \text{and} \quad \langle v, (1, 2, 1) \rangle = 0.
	\]

	These inner products give the following system of equations:

	1. $x_1 + i x_3 = 0$
	2. $x_1 + 2x_2 + x_3 = 0$

	Solving this system, we find:

	- From the first equation, $x_1 = -i x_3$.
	- Substituting into the second equation, we get $-i x_3 + 2x_2 + x_3 = 0$, or $(1 - i) x_3 = -2x_2$.

	Thus, we express $x_2$ and $x_1$ in terms of $x_3$:

	\[
		x_2 = \frac{-(1 - i)}{2} x_3, \quad x_1 = -i x_3.
	\]

	Therefore, the general solution for $v \in S^{\perp}$ is:

	\[
		v = x_3 \left( -i, \frac{-(1 - i)}{2}, 1 \right).
	\]

	Thus, $S^{\perp}$ is spanned by the vector:

	\[
		S^{\perp} = \left\langle \left( -i, \frac{-(1 - i)}{2}, 1 \right) \right\rangle.
	\]
\end{homeworkProblem}

% Homework problem 22 sec 6.2 problem 9
\begin{homeworkProblem}
	Let $W = \text{span}\left( \{(i, 0, 1)\} \right)$ in $\mathbb{C}^3$. Find orthogonal bases for $W$ and $W^{\perp}$.

	\textbf{Solution:}

	The subspace $W$ is spanned by the single vector $w = (i, 0, 1)$. Since $W$ is one-dimensional, an orthogonal basis for $W$ is simply:

	\[
		\{(i, 0, 1)\}.
	\]

	Next, we compute an orthogonal basis for the orthogonal complement $W^{\perp}$. Let $v = (x_1, x_2, x_3) \in W^{\perp}$. For $v$ to be orthogonal to $w$, we require:

	\[
		\langle v, w \rangle = 0,
	\]

	where the inner product in $\mathbb{C}^3$ is defined as $\langle v, w \rangle = x_1 \overline{i} + x_2 \overline{0} + x_3 \overline{1}$. This gives the equation:

	\[
		x_1 (-i) + x_3 = 0.
	\]

	This simplifies to:

	\[
		x_1 = i x_3.
	\]

	Thus, the general solution for $v \in W^{\perp}$ is:

	\[
		v = (i x_3, x_2, x_3).
	\]

	Letting $x_3 = 1$ and $x_2 = 0$, we obtain the vector:

	\[
		v_1 = (i, 0, 1).
	\]

	Thus, an orthogonal basis for $W^{\perp}$ is:

	\[
		W^{\perp} = \{(i, 0, 1)\}.
	\]

	Therefore, we conclude that $W^{\perp}$ has the vector:
	\[
		(0, 0, 0).
	\]
\end{homeworkProblem}

% Homework problem 23 sec 6.2 problem 19a
\begin{homeworkProblem}
	In each of the following parts, find the orthogonal projection of the given vector on the given subspace $W$ of the inner product space
	\[
		V = \mathbb{R}^2, \quad u = (2,6), \quad \text{and} \quad W = \{(x, y): y = 4x\}.
	\]

	\textbf{Solution:}

	The subspace $W$ is given by all vectors of the form $(x, 4x)$, which is a line with slope 4. The orthogonal projection of a vector $u = (2,6)$ onto $W$ can be found by projecting $u$ onto a vector that lies in $W$.

	Let the vector $v = (x, 4x)$ be an arbitrary vector in $W$. The projection of $u$ onto $W$ is the vector $p \in W$ that minimizes the distance between $u$ and $v$. The formula for the projection of a vector $u$ onto a subspace spanned by a vector $v$ is:

	\[
		\text{proj}_W(u) = \frac{\langle u, v \rangle}{\langle v, v \rangle} v.
	\]

	We choose the vector $v = (1, 4)$ because it represents the direction of $W$.

	First, compute the inner products $\langle u, v \rangle$ and $\langle v, v \rangle$:

	\[
		\langle u, v \rangle = \langle (2, 6), (1, 4) \rangle = 2 \cdot 1 + 6 \cdot 4 = 2 + 24 = 26,
	\]

	\[
		\langle v, v \rangle = \langle (1, 4), (1, 4) \rangle = 1^2 + 4^2 = 1 + 16 = 17.
	\]

	Now, apply the projection formula:

	\[
		\text{proj}_W(u) = \frac{26}{17} \cdot (1, 4) = \left( \frac{26}{17}, \frac{104}{17} \right).
	\]

	Thus, the orthogonal projection of $u$ onto $W$ is:

	\[
		\boxed{\left( \frac{26}{17}, \frac{104}{17} \right)}.
	\]
\end{homeworkProblem}

\end{document}
